\documentclass[10pt]{article}
\usepackage[margin=1in]{geometry}
\usepackage[T1]{fontenc}
\usepackage{lmodern,microtype,amsmath,amssymb,amsthm}
\theoremstyle{definition}
\newtheorem{assumption}{Assumption}
\begin{document}
\section*{Chi-Square Constrained Robust Optimization}
We minimize a robust objective $F$ over a closed convex set $C$ in a separable real Hilbert space $E$. Here $F(x)$ is the worst expected sample loss under nonnegative densities relative to a probability law $\nu$, with mean one and chi-square divergence at most $\rho>0$. Sample losses are jointly measurable, convex and bounded on $C$, with jointly measurable bounded subgradients. Each oracle call returns the projection of the query onto $C$, together with its loss and subgradient evaluated using one fresh independent sample from $\nu$. Given $\varepsilon>0$, the goal is a measurable feasible output $\bar{x}$ with $\mathbb{E}[F(\bar{x})-F(x^\star)]\leq\varepsilon$, where $x^\star$ minimizes $F$ on $C$ and expectation includes sampling and algorithm randomness.

Let $\Xi$ be a measurable sample space carrying the probability measure $\nu$. Write $\ell:E\times\Xi\to\mathbb{R}$ for the sample loss and $g:E\times\Xi\to E$ for its supplied subgradient. Define the density class
\[
\mathcal{W}_{\rho}=\left\{
\begin{array}{@{}l@{}}
w:\Xi\to\mathbb{R}\text{ measurable}:\quad
w(\xi)\geq0\text{ for every }\xi,\\
w\in L^2(\nu),\quad \mathbb{E}_{\xi\sim\nu}[w(\xi)]=1,\\
\frac12\mathbb{E}_{\xi\sim\nu}[(w(\xi)-1)^2]\leq\rho
\end{array}
\right\}.
\]
The objective and optimization problem are
\[
F(x)=\sup_{w\in\mathcal{W}_{\rho}}
\mathbb{E}_{\xi\sim\nu}[w(\xi)\ell(x,\xi)],
\qquad
\min_{x\in C}F(x).
\]
The norm $\|\cdot\|$ is induced by the inner product on $E$. The algorithm receives $E$, $\Xi$, $\rho>0$, $G>0$, $B>0$, $R\geq0$, $\varepsilon>0$, and a given point $x^0\in E$; the instance consists of $C$, $\nu$, $\ell$, and $g$, accessed through the oracle described below.

Our assumptions:
\begin{assumption}[Feasible set and projection]
The set $C\subseteq E$ is closed and convex, and $x^0\in C$. There is a measurable map $\Pi_C:E\to C$ satisfying
\[
\langle y-\Pi_C(y),u-\Pi_C(y)\rangle\leq0,
\qquad y\in E,\quad u\in C.
\]
\end{assumption}

\begin{assumption}[Sample losses and subgradients]
The maps $\ell$ and $g$ are jointly measurable. For every $\xi\in\Xi$ and $x,y\in C$,
\[
\begin{aligned}
\ell(x,\xi)+\langle g(x,\xi),y-x\rangle
&\leq\ell(y,\xi),\\
0\leq\ell(x,\xi)&\leq B,\\
\|g(x,\xi)\|&\leq G.
\end{aligned}
\]
\end{assumption}

\begin{assumption}[Initial distance]
There exists $x^\star\in C$ such that
\[
\begin{aligned}
F(x^\star)&\leq F(x) &&\text{for every }x\in C,\\
\|x^0-x^\star\|&\leq R.
\end{aligned}
\]
\end{assumption}

Queries may be made at any $y\in E$. At call $k\geq0$, the oracle returns
\[
\bigl(\Pi_C(y),\ \ell(\Pi_C(y),\xi_k),\
 g(\Pi_C(y),\xi_k)\bigr).
\]
The sample variables $(\xi_k)_{k\geq0}$ are measurable, mutually independent, and each has law $\nu$. The loss and subgradient in a reply use the same sample. Any algorithm randomness is independent of the entire sample sequence.

After the allotted calls, the output $\bar{x}$ must be measurable and belong to $C$ for every realization. Its objective gap must be integrable and satisfy
\[
\mathbb{E}_{P}\bigl[F(\bar{x})-F(x^\star)\bigr]\leq\varepsilon,
\]
where $P$ is the joint law of the samples and the algorithm's randomness.
\end{document}
